%% Lecture 21 -- Retarded Potentials: Physical Interpretation, Far-Field Expansion, and Dipole Radiation
\section{Lecture 21 -- Retarded Potentials, Far-Field Expansion, and Dipole Radiation}\label{sec:lec21}

\begin{center}
\begin{tabular}{ll}
  \textbf{Date:}\quad 16/06/2026 & \\
\end{tabular}
\end{center}
\hrule\vspace{1.5em}

\subsection*{Overview}

Building on the Green-function formalism of Lecture~20, this lecture has three acts.
First, we revisit the physical interpretation of the \emph{retarded} Green function and
explain why the \emph{advanced} solution --- while mathematically valid --- is physically
inadmissible for initial-value problems.
Second, we show that the retarded potentials reduce to Coulomb/Biot--Savart in the static
limit and briefly note the Liénard--Wiechert potentials for a single moving charge.
Third, and most importantly, we derive the radiation fields produced by a \emph{localized},
\emph{slowly-moving} source as seen from a distant observation point. Under two
approximations --- geometric ($r \gg L$) and non-relativistic ($v \ll c$, equivalently
$\lambda \gg L$) --- the retarded potentials expand systematically. The leading
radiation contribution is controlled entirely by the second time derivative of the
\textbf{electric dipole moment} $\bvec{p}(t)$. This is the origin of \textbf{dipole
radiation}.

\paragraph{Notation.}
$\bvec{x}$ is the observation point, $\bvec{y}$ is a source point, $r=|\bvec{x}|$,
$\uvec{r}=\bvec{x}/r$, $R=|\bvec{x}-\bvec{y}|$, and $t_R \equiv t - R/c$ is the
retarded time.

%-----------------------------------------------------------
\subsection{Retarded Potentials: Recap and Physical Interpretation}
\label{subsec:lec21-recap}

\subsubsection*{Physical Motivation}
In Lecture~20 we showed that inserting the retarded Green function
$G_{\rm ret}(x-y)=\delta(c(t-t')-|\bvec{x}-\bvec{y}|)/(4\pi|\bvec{x}-\bvec{y}|)$
into the particular-solution formula and performing the $dt'$ integral via the delta
function yields:

\dfn[def:lec21-ret-potentials]{Retarded Potentials}{
\begin{equation}\label{eq:lec21-phi-ret}
  \phi(t,\bvec{x}) = \int d^3y\,\frac{\rho(t_R,\bvec{y})}{|\bvec{x}-\bvec{y}|},
  \qquad
  \Avec(t,\bvec{x}) = \frac{1}{c}\int d^3y\,\frac{\Jvec(t_R,\bvec{y})}{|\bvec{x}-\bvec{y}|},
\end{equation}
where
\begin{equation}\label{eq:lec21-tret}
  \boxed{t_R \equiv t - \frac{|\bvec{x}-\bvec{y}|}{c}.}
\end{equation}
}

\textit{Significance:} The structure is the same Coulomb/Biot--Savart superposition as in
statics, but the source at $\bvec{y}$ is evaluated not ``now'' but at the \emph{retarded
time} $t_R$: the moment when a light signal had to leave $\bvec{y}$ in order to reach
$\bvec{x}$ exactly at time $t$. Information propagates at $c$, so the field at $\bvec{x}$
now was caused by the source in the past.

\begin{remark}
Because $t_R$ depends on the integration variable $\bvec{y}$, different volume elements
in the source contribute to the field at $(t,\bvec{x})$ at \emph{different} retarded
times. A snapshot of the sources at a single instant does not determine the field at a
distant point; one must account for the varying light-travel time across the source.
\end{remark}

%-----------------------------------------------------------
\subsection{Why the Advanced Green Function Is Unphysical}
\label{subsec:lec21-advanced}

\subsubsection*{Conceptual Background}
The wave equation admits a two-parameter family of Green functions:
\[
  G = a\,\frac{\delta(x^0 - r)}{4\pi r} + b\,\frac{\delta(x^0 + r)}{4\pi r},
  \qquad a + b = \frac{1}{4\pi}.
\]
The $b$-piece, $G_{\rm adv} \propto \delta(x^0+r)/r$, is nonzero only on the
\emph{past} light cone of the observation event: it would make the field at $\bvec{x}$
now depend on sources at future times $t' = t + R/c > t$.

Why can we not use such a solution physically? Consider the equation of motion for a
charge $q$ in the field it generates:
\begin{enumerate}
  \item The force on $q$ at time $t$ requires $\phi$ and $\Avec$ at $(t,\bvec{x}(t))$.
  \item With an advanced contribution, $\phi$ depends on where $q$ will be at $t'>t$.
  \item But the future location $\bvec{x}(t')$ is exactly what the equation of motion
        is supposed to \emph{predict} from the current state.
\end{enumerate}
The argument is circular: solving the equation of motion requires knowing the future,
but the future is what solving the equation of motion is supposed to give.
The retarded Green function avoids this by making the force at $t$ depend only on the
\emph{past} trajectory, which is already known.

\textit{Significance:} Physics selects $b=0$ not by mathematical necessity but by the
requirement that the theory admits initial-value (predictive) formulations. The
choice $b=0$ is the field-theory analogue of specifying initial position and velocity
in Newtonian mechanics.

\nt{Pitfall: the action $S=\int d^4x\,\mathcal{L}$ is invariant under time-reversal, and
so are the Maxwell equations themselves. The \emph{solution} with retarded boundary
conditions spontaneously breaks this symmetry: the field ``remembers'' the past, not
the future. Symmetry of the equations of motion does not imply symmetry of the solutions.}

%-----------------------------------------------------------
\subsection{Static Limit and Liénard--Wiechert Potentials}
\label{subsec:lec21-static-lw}

\subsubsection*{Conceptual Background}
\paragraph{Static limit.}
If $\rho$ and $\Jvec$ are time-independent, then $\rho(t_R,\bvec{y}) = \rho(\bvec{y})$
regardless of $t_R$. The retarded potentials \eqref{eq:lec21-phi-ret} then reduce
immediately to the Coulomb and Biot--Savart integrals from electro- and magnetostatics:
\begin{equation}\label{eq:lec21-static-limit}
  \phi(\bvec{x}) = \int d^3y\,\frac{\rho(\bvec{y})}{|\bvec{x}-\bvec{y}|},
  \qquad
  \Avec(\bvec{x}) = \frac{1}{c}\int d^3y\,\frac{\Jvec(\bvec{y})}{|\bvec{x}-\bvec{y}|}.
\end{equation}
This is a consistency check: retardation disappears when nothing changes.

\paragraph{Liénard--Wiechert potentials (named, not derived).}
For a single point charge $e$ moving along a prescribed trajectory $\bvec{r}(t)$, the
charge and current densities are
\[
  \rho(t,\bvec{y}) = e\,\delta^{(3)}(\bvec{y}-\bvec{r}(t)),
  \qquad
  \Jvec(t,\bvec{y}) = e\,\dot{\bvec{r}}(t)\,\delta^{(3)}(\bvec{y}-\bvec{r}(t)).
\]
Substituting into \eqref{eq:lec21-phi-ret} and performing the integral (with care about
the Jacobian of the retarded-time constraint) yields the
\textbf{Liénard--Wiechert potentials}. We will not derive them explicitly; instead we
proceed to the far-field approximation, which is more transparent and covers the most
physically important regime.

\nt{A key consequence of the Liénard--Wiechert potentials is that a charge in
\emph{uniform motion} does not radiate. Only \emph{accelerating} charges produce
fields that carry energy to spatial infinity.  We will re-derive this from the
far-field expansion below.}

%-----------------------------------------------------------
\subsection{Far-Field Setup: Two Approximations}
\label{subsec:lec21-farfield-setup}

\subsubsection*{Physical Motivation}
We consider a source confined to a spatial region of characteristic size $L$: the
sources $\rho(t,\bvec{y})$ and $\Jvec(t,\bvec{y})$ are negligible for
$|\bvec{y}|\gtrsim L$. The observation point $\bvec{x}$ is far away:
\begin{equation}\label{eq:lec21-farfield-geo}
  r = |\bvec{x}| \gg L.
\end{equation}
This covers the physically ubiquitous case of a radio antenna, an oscillating atom, or
any compact radiating system viewed from a large distance.

\begin{figure}[h]
\centering
\begin{tikzpicture}[>=Stealth, scale=1.2]
  \draw[fill=blue!10, draw=blue!50!black, thick, rounded corners=6pt]
    (-0.65,-0.45) rectangle (0.65,0.45);
  \node at (0,0) {\small sources};
  \draw[<->, thick] (-0.65,-0.65) -- (0.65,-0.65) node[midway, below] {${\sim}L$};
  \node[circle, fill=red!70!black, inner sep=2.5pt] (X) at (5.5,1.0) {};
  \node[right] at (5.5,1.0) {$\bvec{x}$};
  \draw[->, thick, black] (0,0) -- (5.5,1.0) node[midway, above left] {$r \gg L$};
  \draw[->, thick, blue!70!black] (0,0) -- (0.35,0.25) node[right] {$\bvec{y}$};
  \draw[->, thick, green!55!black] (0.35,0.25) -- (5.5,1.0)
    node[midway, below right] {$\bvec{x}-\bvec{y}$};
\end{tikzpicture}
\caption{Far-field geometry. The source region has size ${\sim}L$. The observation
point $\bvec{x}$ satisfies $r=|\bvec{x}|\gg L$. Within the source, $|\bvec{y}|\lesssim L\ll r$.}
\label{fig:lec21-farfield}
\end{figure}

\subsubsection*{Mathematical Setup: Geometric Approximation}
For $|\bvec{y}|\ll r$ we expand
\begin{align}\label{eq:lec21-R-expand}
  R \equiv |\bvec{x}-\bvec{y}|
  &= \sqrt{r^2 - 2\bvec{x}\cdot\bvec{y}+|\bvec{y}|^2}
   = r\sqrt{1 - \frac{2\,\uvec{r}\cdot\bvec{y}}{r}+\frac{|\bvec{y}|^2}{r^2}} \notag\\
  &\approx r\!\left(1 - \frac{\uvec{r}\cdot\bvec{y}}{r}\right)
   = r - \uvec{r}\cdot\bvec{y} + O(L^2/r).
\end{align}
Therefore, at leading order in $L/r$:
\begin{equation}\label{eq:lec21-Rleading}
  \boxed{\frac{1}{R} \approx \frac{1}{r},
  \qquad
  t_R(\bvec{y}) \approx \underbrace{\left(t-\frac{r}{c}\right)}_{t_R^{(0)}}
                         + \frac{\uvec{r}\cdot\bvec{y}}{c}.}
\end{equation}
\textit{Significance:} The $1/R$ prefactor in the integrals \eqref{eq:lec21-phi-ret}
is replaced by the constant $1/r$ (valid because corrections are $O(L/r)$ smaller). The
variation of $R$ with $\bvec{y}$ matters, however, inside the retarded-time argument,
because it produces a phase shift $\uvec{r}\cdot\bvec{y}/c$ that drives the multipole
expansion.

\subsubsection*{Non-Relativistic Approximation and the Taylor Expansion}
The phase shift $\uvec{r}\cdot\bvec{y}/c \lesssim L/c$ is the extra retardation due to
the finite size of the source. For the Taylor expansion of $\rho(t_R(\bvec{y}),\bvec{y})$
around $t_R^{(0)}$ to be valid, this shift must be small compared to the characteristic
time $T$ over which $\rho$ changes significantly:
\begin{equation}\label{eq:lec21-nonrel-condition}
  \frac{L}{c} \ll T \quad\Longleftrightarrow\quad
  v_{\rm typ} = \frac{L}{T} \ll c.
\end{equation}
In words: the typical velocity of charges in the source must be non-relativistic.
Equivalently, since $\lambda = cT$ is the wavelength of the emitted radiation,
\begin{equation}\label{eq:lec21-wavelength}
  \boxed{\lambda = cT \gg L.}
\end{equation}
\textit{Significance:} This is the \emph{long-wavelength} or \emph{electrically small}
approximation: the source fits many times inside one wavelength. Under this condition
we may Taylor-expand the source in the small parameter $\uvec{r}\cdot\bvec{y}/c$:
\begin{equation}\label{eq:lec21-taylor-rho}
  \rho(t_R(\bvec{y}),\bvec{y})
  \approx \rho(t_R^{(0)},\bvec{y})
  + \dot\rho(t_R^{(0)},\bvec{y})\frac{\uvec{r}\cdot\bvec{y}}{c}
  + \ldots
\end{equation}
and similarly for $\Jvec$.

\nt{For an atom of size $L \sim 1\,\text{\AA}$ emitting visible light
$\lambda\sim 500\,\text{nm}$, the ratio $L/\lambda \sim 10^{-3}$. The dipole
approximation is excellent for atomic radiation. For a radio antenna of length
$L \sim 1\,\text{m}$ and frequency $f=300\,\text{MHz}$ ($\lambda=1\,\text{m}$),
we have $L/\lambda \sim 1$ and the approximation breaks down --- higher multipoles
are needed.}

%-----------------------------------------------------------
\subsection{Far-Field Electric Potential}
\label{subsec:lec21-farfield-phi}

\subsubsection*{Mathematical Setup}
Applying \eqref{eq:lec21-Rleading} to the scalar potential in \eqref{eq:lec21-phi-ret}:
\begin{equation}\label{eq:lec21-phi-step1}
  \phi(t,\bvec{x}) \approx \frac{1}{r}
  \int d^3y\,\rho\!\left(t_R^{(0)}+\frac{\uvec{r}\cdot\bvec{y}}{c},\,\bvec{y}\right).
\end{equation}
Inserting the Taylor expansion \eqref{eq:lec21-taylor-rho}:
\begin{equation}\label{eq:lec21-phi-taylor}
  \phi(t,\bvec{x}) \approx
  \frac{1}{r}\underbrace{\int d^3y\,\rho(t_R^{(0)},\bvec{y})}_{\text{total charge }Q}
  + \frac{1}{cr}\underbrace{\int d^3y\,\dot\rho(t_R^{(0)},\bvec{y})\,(\uvec{r}\cdot\bvec{y})}_{\text{dipole term}}
  + \ldots
\end{equation}

\paragraph{First term: total charge.}
Since $Q = \int d^3y\,\rho$ is conserved and time-independent, this gives $Q/r$,
the familiar Coulomb monopole.

\paragraph{Second term: dipole moment.}
Because $t_R^{(0)}$ is independent of $\bvec{y}$, the time derivative can be pulled
out of the integral:
\[
  \int d^3y\,\dot\rho(t_R^{(0)},\bvec{y})\,\bvec{y}
  = \frac{d}{dt}\bigg|_{t=t_R^{(0)}}\int d^3y\,\rho(t,\bvec{y})\,\bvec{y}
  = \dot{\bvec{p}}(t_R^{(0)}),
\]
where $\bvec{p}(t) \equiv \int d^3y\,\rho(t,\bvec{y})\,\bvec{y}$ is the electric dipole
moment. Hence:
\begin{equation}\label{eq:lec21-phi-far}
  \boxed{\phi(t,\bvec{x}) \approx \frac{Q}{r}
  + \frac{\uvec{r}\cdot\dot{\bvec{p}}(t_R^{(0)})}{cr}
  + O\!\left(\frac{L}{r}\right).}
\end{equation}
\textit{Significance:} The scalar potential has a static Coulomb piece plus a dynamical
piece proportional to $\dot{\bvec{p}}$. Both decay as $1/r$, but only the
$\dot{\bvec{p}}$ term is sensitive to the motion of the sources. For a neutral system
($Q=0$), the leading term is the dipole contribution.

\nt{Pitfall: both the Coulomb term $Q/r$ and the dipole correction $\dot{p}/cr$ decay
as $1/r$, yet only the dipole term contributes to radiation. The reason is that
radiation requires the Poynting vector $\bvec{S} \propto \bvec{E}\times\bvec{B}$ to
give a finite total power through a sphere at $r\to\infty$. The Coulomb field gives
$\bvec{E}\sim Q/r^2$, so $|\bvec{S}|\sim Q^2/r^4$, and the total power
$\sim r^2|\bvec{S}|\sim Q^2/r^2\to 0$. Only the fields that decay as $1/r$ in the
\emph{fields} (not the potentials) carry energy to infinity.}

%-----------------------------------------------------------
\subsection{Far-Field Vector Potential}
\label{subsec:lec21-farfield-A}

\subsubsection*{Mathematical Setup}
The same geometric approximation gives:
\begin{equation}\label{eq:lec21-A-step1}
  \Avec(t,\bvec{x}) \approx \frac{1}{cr}
  \int d^3y\,\Jvec\!\left(t_R^{(0)}+\frac{\uvec{r}\cdot\bvec{y}}{c},\,\bvec{y}\right).
\end{equation}
At zeroth order in the Taylor expansion (the leading non-trivial term for $\Avec$):
\begin{equation}\label{eq:lec21-A-step2}
  \Avec(t,\bvec{x}) \approx \frac{1}{cr}\int d^3y\,\Jvec(t_R^{(0)},\bvec{y}).
\end{equation}
We convert the current integral to a dipole derivative using the continuity equation
$\partial_t\rho + \divg\Jvec = 0$. Multiply by $y_i$ and integrate over all space;
integrating the second term by parts (localized source $\Rightarrow$ no boundary term):
\[
  0 = \partial_t\!\int d^3y\,\rho\,y_i + \int d^3y\,y_i\,\nabla_j J_j
    = \dot{p}_i - \int d^3y\,J_i.
\]
Therefore:
\begin{equation}\label{eq:lec21-dpdt-J}
  \boxed{\dot{\bvec{p}}(t) = \int d^3y\,\Jvec(t,\bvec{y}).}
\end{equation}
\textit{Significance:} The rate of change of the dipole moment equals the total current.
This identity (a consequence of charge conservation) converts a vector current integral
into the time derivative of a scalar-weighted charge distribution.

Substituting into \eqref{eq:lec21-A-step2}:
\begin{equation}\label{eq:lec21-A-far}
  \boxed{\Avec(t,\bvec{x}) \approx \frac{\dot{\bvec{p}}(t_R^{(0)})}{cr}.}
\end{equation}
\textit{Significance:} The far-field vector potential is controlled entirely by
$\dot{\bvec{p}}$ at the retarded time. Together with \eqref{eq:lec21-phi-far}, we
see that both the correction to $\phi$ and $\Avec$ are $O(v/c)$ relative to the
Coulomb term and are of the same order in the approximation --- keeping one without
the other would be inconsistent.

%-----------------------------------------------------------
\subsection{Magnetic Field in the Radiation Zone}
\label{subsec:lec21-B-field}

\subsubsection*{Mathematical Setup}
We compute $\bvec{B} = \nabla\times\Avec$ using \eqref{eq:lec21-A-far}:
\[
  \bvec{B} = \nabla\times\!\left[\frac{\dot{\bvec{p}}(t_R^{(0)})}{cr}\right]
  = \frac{1}{cr}\nabla\times\dot{\bvec{p}}(t_R^{(0)})
    + \nabla\!\left(\frac{1}{cr}\right)\times\dot{\bvec{p}}(t_R^{(0)}).
\]
The second term involves $\nabla(1/r) \sim 1/r^2$, giving a $B\sim 1/r^2$
contribution. Since radiation requires fields decaying as $1/r$, we discard this term.

For the first term, the gradient acts through the retarded time
$t_R^{(0)} = t - r/c$:
\[
  [\nabla\times\dot{\bvec{p}}(t_R^{(0)})]_i
  = \epsilon_{ijk}\,\partial_j\dot{p}_k(t - r/c)
  = \epsilon_{ijk}\,\ddot{p}_k(t_R^{(0)})\,\partial_j\!\left(-\frac{r}{c}\right)
  = -\frac{\epsilon_{ijk}\,\hat{r}_j\,\ddot{p}_k}{c}.
\]
Recognising the cross product $(\bvec{A}\times\bvec{B})_i = \epsilon_{ijk}A_jB_k$:
\[
  \nabla\times\dot{\bvec{p}}(t_R^{(0)}) = -\frac{\ddot{\bvec{p}}(t_R^{(0)})\times\uvec{r}}{c}.
\]
Therefore:
\begin{equation}\label{eq:lec21-B-rad}
  \boxed{\bvec{B}(t,\bvec{x}) = \frac{\ddot{\bvec{p}}(t_R^{(0)})\times\uvec{r}}{c^2 r}.}
\end{equation}
\textit{Significance:} The magnetic radiation field is proportional to $\ddot{\bvec{p}}$
--- the \emph{second} time derivative of the dipole moment, i.e.\ the acceleration of
the charges. A dipole oscillating at constant velocity ($\ddot{\bvec{p}}=0$) produces no
radiation magnetic field. The factor $1/r$ ensures a finite energy flux through any
sphere.

\nt{Pitfall: when computing $\nabla\times\Avec$ in the radiation zone, \emph{do not}
differentiate the $1/r$ prefactor --- that yields $1/r^2$ and belongs to the near
field. The radiation arises entirely from differentiating \emph{through} the retarded
time $t_R^{(0)}=t-r/c$, which brings down one factor of $\ddot{\bvec{p}}/c$.}

%-----------------------------------------------------------
\subsection{Electric Field in the Radiation Zone}
\label{subsec:lec21-E-field}

\subsubsection*{Mathematical Setup}
The electric field is $\bvec{E}=-\nabla\phi - (1/c)\partial_t\Avec$.
Differentiating $\Avec = \dot{\bvec{p}}(t_R^{(0)})/(cr)$ with respect to $t$ (at
fixed $\bvec{x}$) gives $(1/c)\partial_t\Avec = \ddot{\bvec{p}}(t_R^{(0)})/(c^2r)$.
Differentiating the dipole term in $\phi$ (again only through $t_R^{(0)}$) gives
$-\nabla\phi_{\rm dip} \approx +\uvec{r}(\uvec{r}\cdot\ddot{\bvec{p}})/(c^2r)$.
Adding these two contributions (the derivation is left as an exercise and is analogous
to the $\bvec{B}$ calculation):

\begin{equation}\label{eq:lec21-E-rad}
  \boxed{\bvec{E}(t,\bvec{x})
  = \frac{[\ddot{\bvec{p}}(t_R^{(0)})\times\uvec{r}]\times\uvec{r}}{c^2 r}.}
\end{equation}
\textit{Significance:} The electric radiation field is the part of $\ddot{\bvec{p}}$
transverse to the line of sight. Expanding with the BAC--CAB rule,
$(\ddot{\bvec{p}}\times\uvec{r})\times\uvec{r}
= \uvec{r}(\uvec{r}\cdot\ddot{\bvec{p}}) - \ddot{\bvec{p}}
= -\ddot{\bvec{p}}_\perp$,
so $\bvec{E}=-\ddot{\bvec{p}}_\perp/(c^2r)$: the field vanishes along the dipole axis
($\theta=0$) and is maximum in the equatorial plane ($\theta=\pi/2$).

\subsubsection*{Plane-Wave Structure and the Radiation Triad}
A direct computation using \eqref{eq:lec21-B-rad} and \eqref{eq:lec21-E-rad} gives:
\begin{equation}\label{eq:lec21-E-from-B}
  \boxed{\bvec{E} = \bvec{B}\times\uvec{r},}
\end{equation}
which implies that $\bvec{E}$, $\bvec{B}$, and $\uvec{r}$ form a right-handed orthogonal
triad with $|\bvec{E}|=|\bvec{B}|$. This is exactly the structure of a plane wave
propagating in the $\uvec{r}$ direction.

\textit{Significance:} At large distances a spherical wave looks locally flat, just as
the curved surface of the Earth appears flat on human scales. The energy flows radially
outward via the Poynting vector $\bvec{S} = (c/4\pi)\bvec{E}\times\bvec{B}$,
which points in the $\uvec{r}$ direction.

\begin{figure}[h]
\centering
\begin{tikzpicture}[>=Stealth, scale=1.4]
  % Dipole axis
  \draw[->, very thick, black!70] (0,-0.8) -- (0,0.8)
    node[right] {$\ddot{\bvec{p}}$};
  % r direction
  \draw[->, thick, red!70!black] (0,0) -- (2.8,1.6)
    node[right] {$\uvec{r}$};
  % theta arc
  \draw[gray] (0,0.45) arc (90:30:0.45) node[midway, right] {\small$\theta$};
  % B field (perpendicular to both r and p, into/out of page)
  \draw[->, thick, blue!70!black] (2.8,1.6) -- (1.8,2.5)
    node[above] {$\bvec{B}$};
  % E field
  \draw[->, thick, green!60!black] (2.8,1.6) -- (3.8,1.2)
    node[right] {$\bvec{E}$};
  % No radiation along the axis
  \draw[dashed, gray!60] (0,0) -- (0,2.2);
  \node[gray, right] at (0.08,2.0) {\small no radiation ($\theta=0$)};
  % Maximum radiation label
  \node[black!60, below right] at (2.2, -0.2) {\small max radiation ($\theta=\pi/2$)};
  \draw[dashed, gray!60] (0,0) -- (2.5,-0.3);
\end{tikzpicture}
\caption{Radiation fields of an oscillating electric dipole. Both $\bvec{E}$ and
$\bvec{B}$ are perpendicular to $\uvec{r}$ and to each other. The intensity is
proportional to $\sin^2\theta$: maximum in the equatorial plane, zero along the
dipole axis.}
\label{fig:lec21-dipole-pattern}
\end{figure}

%-----------------------------------------------------------
\subsection{Dipole Approximation: Summary and Scope}
\label{subsec:lec21-dipole-summary}

\dfn[def:lec21-dipole-approx]{Electric Dipole (Long-Wavelength) Approximation}{
Valid when $r \gg L$ (far field) and $v\ll c$ (equivalently $\lambda \gg L$).
The leading radiation fields are, with $t_R = t - r/c$:
\begin{align*}
  \Avec &= \frac{\dot{\bvec{p}}(t_R)}{cr},
  &
  \bvec{B} &= \frac{\ddot{\bvec{p}}(t_R)\times\uvec{r}}{c^2r},
  &
  \bvec{E} &= \bvec{B}\times\uvec{r},
\end{align*}
where $\bvec{p}(t)=\int d^3y\,\rho(t,\bvec{y})\,\bvec{y}$ is the electric dipole
moment of the source.
}

\subsubsection*{Method: Computing Dipole Radiation}
\begin{enumerate}
  \item Confirm the approximation is valid: $r \gg L$ and $v \ll c$ (or $\lambda \gg L$).
  \item Compute $\bvec{p}(t) = \sum_\alpha q_\alpha \bvec{r}_\alpha(t)$ (discrete) or
        $\int d^3y\,\rho\,\bvec{y}$ (continuous).
  \item Differentiate twice to get $\ddot{\bvec{p}}(t)$.
  \item Write $\bvec{B} = (\ddot{\bvec{p}}(t_R)\times\uvec{r})/(c^2r)$ and
        $\bvec{E} = \bvec{B}\times\uvec{r}$.
  \item The angular power distribution follows $dP/d\Omega\propto|\ddot{\bvec{p}}|^2\sin^2\theta$
        (to be derived next lecture).
\end{enumerate}

\ex[ex:lec21-oscillating]{Worked Example: Oscillating Dipole}{
A charge $+q$ oscillates along the $z$-axis: $z(t) = d\cos(\omega t)$, with $-q$ fixed
at the origin. The dipole moment is $\bvec{p}(t) = qd\cos(\omega t)\,\uvec{z}$, so
$\ddot{\bvec{p}}(t) = -qd\omega^2\cos(\omega t)\,\uvec{z}$.

In spherical coordinates ($\uvec{r}$ at polar angle $\theta$ from $z$-axis):
\begin{align*}
  \bvec{B}
  &= \frac{-qd\omega^2\cos(\omega t_R)}{c^2 r}\,\uvec{z}\times\uvec{r}
   = \frac{qd\omega^2\cos(\omega t_R)\sin\theta}{c^2 r}\,\uvec{\phi},\\
  \bvec{E}
  &= \bvec{B}\times\uvec{r}
   = \frac{qd\omega^2\cos(\omega t_R)\sin\theta}{c^2 r}\,\uvec{\theta},
\end{align*}
where $t_R = t - r/c$.

Both fields are maximum in the equatorial plane ($\theta=\pi/2$) and vanish along the
$z$-axis ($\theta = 0,\pi$). The amplitude falls off as $1/r$, confirming that this
field carries energy to infinity: it is genuine radiation.
}

\nt{Pitfall: for a system of charges with total charge $Q = \sum_\alpha q_\alpha \neq 0$,
the centre-of-mass contribution to $\ddot{\bvec{p}}$ is $Q\ddot{\bvec{r}}_{\rm cm}$.
If all charges accelerate identically (rigid-body motion), there is no relative motion
and hence no radiation, even though $\ddot{\bvec{p}} \neq 0$ superficially. A truly
neutral system ($Q=0$) cannot radiate via the centre-of-mass mode; radiation requires
\emph{internal} relative acceleration.}

\subsubsection*{Beyond Dipole: Higher Multipoles}
The Taylor expansion \eqref{eq:lec21-taylor-rho} is a systematic series in $L/\lambda$.
When the dipole contribution vanishes (by symmetry or selection rules), the next
non-zero term in the expansion gives either the \textbf{electric quadrupole} or the
\textbf{magnetic dipole} radiation, both of order $(L/\lambda)^2$ smaller than the
electric dipole. We will not derive these here, but the pattern is clear: higher
multipoles enter at higher order in the long-wavelength expansion.

%-----------------------------------------------------------
\subsection*{Summary}

\begin{itemize}
  \item The retarded potentials \eqref{eq:lec21-phi-ret} express the field at $(t,\bvec{x})$
        as a Coulomb/Biot--Savart superposition evaluated at retarded times
        $t_R = t - |\bvec{x}-\bvec{y}|/c$. Different volume elements contribute at different
        past moments.
  \item The advanced Green function is mathematically valid but physically inadmissible: using
        it would require knowledge of the future to solve the equations of motion, making
        initial-value predictions impossible.
  \item In the static limit the retarded potentials reduce to Coulomb/Biot--Savart.
  \item Under the two approximations $r\gg L$ and $v\ll c$ ($\Leftrightarrow \lambda\gg L$):
    \[
      \phi \approx \frac{Q}{r} + \frac{\uvec{r}\cdot\dot{\bvec{p}}(t_R)}{cr},
      \qquad
      \Avec \approx \frac{\dot{\bvec{p}}(t_R)}{cr}.
    \]
  \item The identity $\dot{\bvec{p}} = \int d^3y\,\Jvec$ (from continuity) converts the
        current integral to a dipole derivative.
  \item The radiation fields decay as $1/r$ and are:
    \[
      \bvec{B} = \frac{\ddot{\bvec{p}}(t_R)\times\uvec{r}}{c^2r},
      \qquad
      \bvec{E} = \bvec{B}\times\uvec{r}.
    \]
  \item Radiation requires $\ddot{\bvec{p}}\neq 0$: only \emph{accelerating} charges
        radiate at dipole order.
  \item At large $r$, $\bvec{E}\perp\bvec{B}\perp\uvec{r}$ with $|\bvec{E}|=|\bvec{B}|$:
        the radiation locally resembles a plane wave.
\end{itemize}

\subsection*{Further Reading}
\begin{itemize}
  \item J.\,D. Jackson, \textit{Classical Electrodynamics}, 3rd ed.: Chapter 9
        (multipole radiation, dipole radiation).
  \item D.\,J. Griffiths, \textit{Introduction to Electrodynamics}, 4th ed.: Chapter 11
        (radiation).
  \item L.\,D. Landau \& E.\,M. Lifshitz, \textit{The Classical Theory of Fields}, Vol.\ 2:
        Chapter 9 (radiation of electromagnetic waves).
  \item David Tong, \textit{Lectures on Electromagnetism}: Chapter 6, radiation.
\end{itemize}
